% docs/texte/05_theory/visibility.tex -- theory chapter "Visibility and Source Models" (Interferosim)
% Version: 0.1.0 (2026-10-12)
% Notation as defined in texte/05_theory/notation.tex.
% ================================================================

\label{ch:visibility}
\index{visibility}\index{van Cittert--Zernike theorem}
\section{The visibility function}

For a field small enough that the $w$ term can be neglected, the
response of a baseline $(u,v)$ to the brightness distribution $I(l,m)$
is the complex visibility
\begin{equation}
  \mathcal V(u,v) = \iint A(l,m)\,I(l,m)\,
     e^{-2\pi i\,(ul+vm)}\,\mathrm{d}l\,\mathrm{d}m .
  \label{eq:vis}
\end{equation}
This is the van Cittert--Zernike theorem: the visibility is the
two-dimensional Fourier transform of the sky brightness weighted by
the primary beam $A(l,m)$. Since $I$ is real,
$\mathcal V(-u,-v) = \mathcal V^*(u,v)$; every baseline therefore
measures two points of the $(u,v)$ plane.

\section{Analytic visibilities}

Interferosim computes the visibilities of its source components
analytically (\texttt{SkyComponent.visibility}). With $q=\sqrt{u^2+v^2}$:

\paragraph*{Point source} of flux density $S$ at $(l_0,m_0)$:
\begin{equation}
  \mathcal V = S\,e^{-2\pi i\,(ul_0+vm_0)} .
\end{equation}

\paragraph*{Gaussian source} with FWHM $\theta_{1/2}$ (circular):
\begin{equation}
  \mathcal V = S\,\exp\!\left(-\frac{\pi^2\,\theta_{1/2}^2\,q^2}{4\ln 2}\right)
  e^{-2\pi i\,(ul_0+vm_0)} .
\end{equation}
For an elliptical source, $\theta_{1/2}^2 q^2$ is replaced by
$\theta_{\mathrm{maj}}^2 u'^2 + \theta_{\mathrm{min}}^2 v'^2$ with
$(u',v')$ rotated by the position angle.

\paragraph*{Uniform disc} of diameter $\theta_d$:
\begin{equation}
  \mathcal V = S\,\frac{2J_1(\pi\theta_d q)}{\pi\theta_d q} .
  \label{eq:vis-disc}
\end{equation}
The first zero lies at $\pi\theta_d q = 3.8317$, i.\,e.\ at
\begin{equation}
  b_{\mathrm{null}} = 1.2197\,\frac{\lambda}{\theta_d} .
\end{equation}
For the Sun at 21\,cm ($\theta_d = 34'$) this is $b_{\mathrm{null}}
\approx 26\,\mathrm{m}$; measuring the position of the null is a
measurement of the radio diameter (notebook~B).

\paragraph*{Circularly symmetric source} with radial profile $I(r)$,
e.\,g.\ a limb-brightened solar disc: the two-dimensional Fourier
transform reduces to a Hankel transform,
\begin{equation}
  \mathcal V(q) = 2\pi\int_0^{R} I(r)\,J_0(2\pi q r)\,r\,\mathrm{d}r .
\end{equation}
Fitting the uniform-disc model \eqref{eq:vis-disc} to data of a
limb-brightened disc biases the diameter: the formal fit error then
underestimates the true error by a large factor, because the model, not
the noise, dominates the error budget.
